% Part: lambda-calculus
% Chapter: syntax
% Section: substitution

\documentclass[../../../include/open-logic-section]{subfiles}

\begin{document}

\olfileid{lam}{syn}{sub}
\olsection{جای‌گذاری}

\begin{explain}
متغیرهای آزاد ارجاع‌هایی به متغیرهای محیط‌اند؛ ازاین‌رو معقول است که
به‌جای یک متغیر آزاد، ارزشی معین را واقعاً به کار ببریم. برای مثال،
ممکن است بخواهیم در~$\lambd[x][f x]$، به‌جای~$f$ ترم معینی، مانند تابع
همانی~$\lambd[y][y]$، قرار دهیم. حاصل~$\lambd[x][(\lambd[y][y]) x]$
است. فرایند جایگزین‌کردن متغیرهای آزاد با ترم‌های لامبدا
جای‌گذاری نام دارد.
\end{explain}

\begin{defn}[جای‌گذاری] \ollabel{defn:substitution}
  \emph{جای‌گذاری} ترم~$N$ به‌جای متغیر~$x$ در ترم~$M$، که با
  $\Subst{M}{N}{x}$ نشان داده می‌شود، به‌صورت استقرایی چنین تعریف می‌شود:
  \begin{enumerate}
    \item $\Subst{x}{N}{x} = N$. \ollabel{defn:substitution-1} 
    \item $\Subst{y}{N}{x}  = y$ اگر~$x \neq y$. \ollabel{defn:substitution-2} 
    \item $\Subst{PQ}{N}{x}  = (\Subst{P}{N}{x}) (\Subst{Q}{N}{x})$.
      \ollabel{def:substitution-3}
    \item $\Subst{(\lambd[y][P])}{N}{x} = \lambd[y][\Subst{P}{N}{x}]$،
      اگر~$x \neq y$ و~$y \notin \FV{N}$؛ وگرنه تعریف‌نشده است.
      \ollabel{defn:substitution-4}
  \end{enumerate}
\end{defn}

\begin{explain}
در~\olref{defn:substitution}\olref{defn:substitution-4}، شرط~$x \neq y$
را می‌گذاریم، زیرا نمی‌خواهیم رخدادهای \emph{مقید} متغیر~$x$ در~$M$ را
با~$N$ جایگزین کنیم. برای مثال، در محاسبهٔ جای‌گذاری
$\Subst{\lambd[x][x]}{y}{x}$، حاصل نباید~$\lambd[x][y]$ باشد. طبق این
تعریفِ عمداً جزئی، این جای‌گذاری تعریف‌نشده است؛ مسئله بعدتر با تبدیل
آلفا و تغییر نام متغیرهای مقید برطرف خواهد شد.

هنگام جای‌گذاری~$N$ به‌جای~$x$ در~$\lambd[y][P]$، همچنین شرط
$y \notin \FV{N}$ را می‌گذاریم. برای مثال، نمی‌توان~$y$ را به‌جای~$x$
در~$\lambd[y][x]$ جای‌گذاری کرد، یعنی
$\Subst{\lambd[y][x]}{y}{x}$، زیرا حاصل آن~$\lambd[y][y]$ می‌شد؛ ترمی
که تابع پذیرندهٔ یک آرگومان و بازگردانندهٔ همان آرگومان را نشان می‌دهد.
اما ترم~$\lambd[y][x]$ تابعی را نشان می‌دهد که همواره ترم~$x$ (یا هرچه
$x$ به آن ارجاع می‌دهد) را بازمی‌گرداند. پس حاصل مطلوب در واقع تابعی
است که آرگومانی می‌پذیرد، آن را کنار می‌گذارد و متغیر محیطی~$y$ را
بازمی‌گرداند. برای انجام درست این کار، نخست باید متغیر مقید~$y$ را
«تغییر نام» دهیم.
\end{explain}

\begin{prob}
  حاصل جای‌گذاری‌های زیر چیست؟
  \begin{enumerate}
  \item $\Subst{\lambd[y][x(\lambd[w][vwx])]}{(uv)}{x}$
  \item $\Subst{\lambd[y][x(\lambd[x][x])]}{(\lambd[y][xy])}{x}$
  \item $\Subst{y(\lambd[v][xv])}{(\lambd[y][vy])}{x}$
  \end{enumerate}
\end{prob}

\begin{thm} \ollabel{thm:notinfv}
  اگر~$x \notin \FV{M}$، آنگاه، مشروط بر آنکه طرف چپ تعریف شده باشد،
  $\FV{\Subst{M}{N}{x}} = \FV{M}$.
\end{thm}

\begin{proof}
  با استقرا بر ساخت~$M$.
  \begin{enumerate}
  \item $M$ یک متغیر است: تمرین.
  \item $M$ به‌صورت~$(PQ)$ است: تمرین.
  \item $M$ به‌صورت~$\lambd[y][P]$ است، و چون
    $\Subst{\lambd[y][P]}{N}{x}$ تعریف شده است، باید برابر با
    $\lambd[y][\Subst{P}{N}{x}]$ باشد. بنابراین~$\Subst{P}{N}{x}$ نیز
    باید تعریف شده باشد؛ افزون بر این، $x \neq y$، $y \notin \FV{N}$
    و~$x \notin \FV{P}$. در نتیجه:
    \begin{multline*}
      \FV{\Subst{\lambd[y][P]}{N}{x}} = \\
    \begin{aligned}
      & = \FV{\lambd[y][\Subst{P}{N}{x}]} &&
      \text{بنا بر \olref{defn:substitution-4}}\\
      & = \FV{\Subst{P}{N}{x}} \setminus \{y\} &&
      \text{بنا بر \olref[fv]{def:fv}\olref[fv]{def:fv2}}\\
      & = \FV{P} \setminus \{y\} && \text{بنا بر فرض استقرا} \\
      & = \FV{\lambd[y][P]} && \text{بنا بر \olref[fv]{def:fv}\olref[fv]{def:fv2}}
    \end{aligned}
    \end{multline*}
  \end{enumerate}
\end{proof}

\begin{prob}
  برهان~\olref[lam][syn][sub]{thm:notinfv} را کامل کنید.
\end{prob}

\begin{thm} \ollabel{thm:infv}
  اگر~$x \in \FV{M}$، آنگاه
  $\FV{\Subst{M}{N}{x}} = (\FV{M} \setminus \{x\}) \cup \FV{N}$،
  مشروط بر آنکه طرف چپ تعریف شده باشد.
\end{thm}

\begin{proof}
  با استقرا بر ساخت~$M$.
  \begin{enumerate}
  \item $M$ یک متغیر است: تمرین.
  \item $M$ به‌صورت~$PQ$ است: چون
    $\Subst{(PQ)}{N}{x}$ تعریف شده است، باید برابر با
    $(\Subst{P}{N}{x})(\Subst{Q}{N}{x})$ باشد و هر دو جای‌گذاری نیز
    تعریف شده باشند. همچنین چون~$x \in \FV{PQ}$، یا~$x \in \FV{P}$،
    یا~$x \in \FV{Q}$، یا هر دو. بقیه به‌عنوان تمرین واگذار می‌شود.
  \item $M$ به‌صورت~$\lambd[y][P]$ است. چون
    $\Subst{\lambd[y][P]}{N}{x}$ تعریف شده است، باید برابر با
    $\lambd[y][\Subst{P}{N}{x}]$ باشد، که در آن~$\Subst{P}{N}{x}$
    تعریف شده، $x \neq y$ و~$y \notin \FV{N}$؛ همچنین چون~$x \in
    \FV{\lambd[y][P]}$، داریم~$x \in \FV{P}$. اکنون:
    \begin{multline*}
      \FV{\Subst{(\lambd[y][P])}{N}{x}} = \\
      \begin{aligned}
      & = \FV{\lambd[y][\Subst{P}{N}{x}]} \\
      & = \FV{\Subst{P}{N}{x}} \setminus \{y\} \\
      &= ((\FV{P}\setminus\{x\})\cup\FV{N})\setminus\{y\}
       && \text{بنا بر فرض استقرا} \\
      &= (\FV{P}\setminus\{x,y\})\cup\FV{N}
       && \text{چون }y\notin\FV{N} \\
      & = (\FV{\lambd[y][P]} \setminus \{x\}) \cup \FV{N}
      \end{aligned}
      \end{multline*}
  \end{enumerate}
\end{proof}

\begin{prob}
  برهان~\olref[lam][syn][sub]{thm:infv} را کامل کنید.
\end{prob}

\begin{thm}\ollabel{thm:clr}
  اگر طرف راست تعریف شده باشد و~$x \notin \FV{N}$، آنگاه
  $x \notin \FV{\Subst{M}{N}{x}}$.
\end{thm}

\begin{proof}
  تمرین.
\end{proof}

\begin{prob}
  \olref[lam][syn][sub]{thm:clr} را ثابت کنید.
\end{prob}


\begin{thm}\ollabel{thm:inv}
  اگر~$\Subst{M}{y}{x}$ تعریف شده باشد و~$y \notin \FV{M}$، آنگاه
  $\Subst{\Subst{M}{y}{x}}{x}{y} = M$.
\end{thm}

\begin{proof}
  با استقرا بر ساخت~$M$.
  \begin{enumerate}
  \item $M$ متغیر~$z$ است: تمرین.
  \item $M$ به‌صورت~$(PQ)$ است. آنگاه:
    \begin{align*}
      \Subst{\Subst{(PQ)}{y}{x}}{x}{y}
      &=\Subst{((\Subst{P}{y}{x})(\Subst{Q}{y}{x}))}{x}{y} \\
      &= (\Subst{\Subst{P}{y}{x}}{x}{y})(\Subst{\Subst{Q}{y}{x}}{x}{y}) \\
      &= (PQ) \text{ بنا بر فرض استقرا}
    \end{align*}
  \item $M$ به‌صورت~$\lambd[z][N]$ است. چون
    $\Subst{\lambd[z][N]}{y}{x}$ تعریف شده است، می‌دانیم
    $z \neq y$. پس:
    \begin{align*}
      \Subst{\Subst{(\lambd[z][N])}{y}{x}}{x}{y}\\
      & = \Subst{(\lambd[z][\Subst{N}{y}{x}])}{x}{y} \\
      & = \lambd[z][\Subst{\Subst{N}{y}{x}}{x}{y}] \\
      &= \lambd[z][N] \text{ بنا بر فرض استقرا} 
    \end{align*}
  \end{enumerate}
\end{proof}

\begin{prob}
  برهان~\olref[lam][syn][sub]{thm:inv} را کامل کنید.
\end{prob}

\end{document}
